In \(\mathfrak G_\varepsilon\), \(p_t=1/2\) for both times.  For a binary nonfactual row whose observed marginal is \((1/2,1/2)\), \(\texttt{def:posterior-row-fiber}\) imposes
\[
 \frac12q(0)\le\frac12,
 \qquad
 \frac12q(1)\le\frac12,
 \qquad
 q(0)+q(1)=1,
\]
which is exactly \(0\le q(1)\le1\).  Because the fiber is a Cartesian product of row factors, \(u,w,z,v\) range independently over \([0,1]\), and the factual rows equal \(\delta_0\).

The first-stage reward is zero and the second-stage reward is the terminal next state.  Under the empty coalition both factual rows are queried, so \(g_\varnothing=0\).  Under \(\{1\}\), the time-one baseline chooses action one with probability \(\varepsilon\), reaches state one with conditional probability \(u\), and the retained time-two policy then uses action zero, whose success probability is \(z\).  Hence
\[
 g_{\{1\}}=\varepsilon uz.
\]
Under \(\{2\}\), time one remains factual and time two takes action one at state zero, giving
\[
 g_{\{2\}}=w.
\]
Under \(\{1,2\}\), the time-one baseline reaches state one with probability \(\varepsilon u\).  The time-two baseline success probability is \(w\) from state zero and \(v\) from state one.  Therefore
\[
 g_{\{1,2\}}=(1-\varepsilon u)w+\varepsilon uv.
\]
For \(H=2\), both Shapley weights are \(1/2\).  Substitution into \(\texttt{def:shapley-coordinate}\) gives
\[
\begin{aligned}
 \Phi_1(q)
 &=\frac12\{g_\varnothing-g_{\{1\}}+g_{\{2\}}-g_{\{1,2\}}\}\\
 &=\frac{\varepsilon u}{2}(w-z-v).
\end{aligned}
\]
Since \(u,w,z,v\in[0,1]\), the maximum is \(\varepsilon/2\), attained at \((u,w,z,v)=(1,1,0,0)\), and the minimum is \(-\varepsilon\), attained at \((1,0,1,1)\).  The equality with \([L_1,U_1]\), together with compatible-SCM attainment by these common row assignments, follows from \(\texttt{thm:sharp-shared-interval}\).

The internally defined relaxation assigns a separate posterior-kernel copy to each distinct coalition return.  The displayed recursions imply
\[
 g_\varnothing\in\{0\},
 \qquad
 g_{\{1\}}\in[0,\varepsilon],
 \qquad
 g_{\{2\}}\in[0,1],
 \qquad
 g_{\{1,2\}}\in[0,1].
\]
Each endpoint is attained by a suitable row assignment in its own copy.  The explicit interval formula in \(\texttt{def:independent-coalition-outer-bound}\) consequently gives
\[
 \mathcal I_1^{\mathrm{ind}}(K,\tau)
 =\left[-\frac{1+\varepsilon}{2},\frac12\right].
\]
Its lower excess below \(-\varepsilon\) and upper excess above \(\varepsilon/2\) both equal
\[
 \frac{1-\varepsilon}{2}>0.
\]

We prove the neighborhood assertion rather than infer it from the single tuple.  On the fixed discrete stratum, restrict first to a sufficiently small relative neighborhood on which both factual transition cells remain positive.  If \(K'\) approaches \(K\) in maximum rowwise \(\ell^1\) distance and their factual cells have a common positive lower bound, truncate each old feasible posterior row at the new caps \(K'_t(y\mid x)/p'_t\) and redistribute the removed mass into unused cap capacity.  Exactly as in the row-matching calculation in \(\texttt{thm:quantitative-gap-neighborhood}\), the removed mass tends uniformly to zero.  Reversing \(K,K'\) proves Hausdorff convergence of the finite product fibers.

Every coalition return is a finite polynomial in the posterior rows, rewards, and policy probabilities, hence is uniformly continuous on a compact neighborhood of the valid primitive and posterior-row products.  Given Hausdorff-close fibers, match a maximizer or minimizer in either direction and apply this uniform continuity.  Thus each shared endpoint and each coalition-return extremum in the internal relaxation is continuous.  The two excess-width functions are finite sums and differences of these continuous extrema.  At \(\varepsilon=1/4\), both equal \(3/8\), so the set on which both are positive is a relative-open neighborhood of the displayed primitive tuple.

This neighborhood genuinely contains interior-policy points.  More explicitly, keep \(K,r\), and the time-one baseline at their \(\varepsilon=1/4\) values, and for \(0\le c\le1/64\) set \(\pi_t(1\mid s)=c\) and \(\pi_{\mathrm{base},2}(1\mid s)=1-c\) for all \(t,s\).  The same recursion, now with \(\eta=c\) and \(\beta=1-c\), gives
\[
 [L_1,U_1]=\left[-\frac14,\frac{1-c}{8}\right],
\]
\[
 \mathcal I_1^{\mathrm{ind}}
 =\left[-\frac58+\frac{3c}{8},\frac{1-c}{2}\right].
\]
The lower and upper excess widths are respectively \(3(1-c)/8\) and \(3(1-c)/8\), so the whole segment remains in the positive-gap set, and every point with \(c>0\) has all policy cells positive.  At \(c=1/64\) this point is \(\theta^\star\), and \(\texttt{thm:quantitative-gap-neighborhood}\) further supplies the explicit radius \(\rho=1/1024\), factual-region mass, support, variance, and uniform \(1/4\)-gap certificate under perturbations of every continuous primitive.  This proves the stated open-neighborhood conclusion and its quantitative strengthening.